\documentclass[ ../main.tex]{subfiles}
\providecommand{\mainx}{..}
\begin{document}
\section{Preservation of relational invariants}
\label{sec:invariants}
Sometimes a relation must satisfy a structural invariant: reflexivity,
symmetry, single-valuedness, transitivity, and so on. An exact relation either
has such a property or it does not. A random approximate relation, by contrast,
tests each pair through an independent Bernoulli trial (\cref{asm:ber_trial}),
and a property that constrains \emph{many} pairs at once is preserved only when
\emph{all} of the relevant trials cooperate. The recurring lesson of this section
is a fragility: the probability that an approximate relation inherits a global
invariant is a product over the constrained pairs, and it tends to zero as the
ground universe grows. Where a cheap repair exists (as for symmetry) we exhibit
it and pay its rate cost; where it does not (as for transitivity) we state the
honest limiting behavior rather than force a clean formula.

Throughout, $\Relation{R} \subseteq \Set{A} \times \Set{A}$ is an exact binary
relation with the named property, and $\ARelation{R}[\fprate][\tprate]$ is its
random approximate relation, with $\tprate = 1 - \fnrate$ the true-positive rate.
We ask, for each property, the probability that $\ARelation{R}$ \emph{also} has it.

\subsection{Reflexivity}
\label{sec:reflexive}
A relation is \emph{reflexive} when it contains the whole diagonal:
$\Pair{a}{a} \in \Relation{R}$ for every $a \in \Set{A}$. The approximation is
reflexive only if each of the $\Card{\Set{A}}$ diagonal pairs survives as a true
positive.

\begin{theorem}[Reflexivity is preserved with vanishing probability]
\label{thm:reflexive}
Let $\Relation{R}$ be reflexive. Then $\ARelation{R}[\fprate][\tprate]$ is
reflexive with probability
\begin{equation}
    \Prob{\ARelation{R} \text{ reflexive}} = \tprate^{\Card{\Set{A}}}\,,
\end{equation}
which tends to $0$ as $\Card{\Set{A}} \to \infty$ for any $\tprate < 1$.
\end{theorem}
\begin{proof}
Each diagonal pair $\Pair{a}{a}$ is a true positive of $\ARelation{R}$ with
probability $\tprate$, and the $\Card{\Set{A}}$ diagonal trials are independent
(\cref{asm:ber_trial}). The approximation is reflexive exactly when all of them
fire, an event of probability $\tprate^{\Card{\Set{A}}}$. Since $\tprate < 1$,
this product decays geometrically in $\Card{\Set{A}}$.
\end{proof}

\subsection{Symmetry}
\label{sec:symmetric}
A relation is \emph{symmetric} when $\Pair{a}{b} \in \Relation{R}$ implies
$\Pair{b}{a} \in \Relation{R}$. In the approximation, the two ordered trials of
an unordered pair are tested independently and need not agree, so symmetry is a
constraint that each unordered pair must satisfy on its own.

\begin{theorem}[Symmetry is preserved with vanishing probability]
\label{thm:symmetric_preservation}
Let $\Relation{R}$ be symmetric, and let $P$ and $N$ be the numbers of unordered
positive and negative pairs (off-diagonal pairs in and out of $\Relation{R}$,
respectively). Then $\ARelation{R}[\fprate][\tprate]$ is symmetric with
probability
\begin{equation}
    \Prob{\ARelation{R} \text{ symmetric}}
    = \left(\tprate^2 + \fnrate^2\right)^{P}
      \left((1-\fprate)^2 + \fprate^2\right)^{N}\,,
\end{equation}
which tends to $0$ as the number of pairs grows, since each base lies strictly
below $1$ whenever $0 < \fprate < 1$ and $0 < \fnrate < 1$.
\end{theorem}
\begin{proof}
Fix an unordered pair $\{a,b\}$ with $a \neq b$. Its two ordered trials, on
$\Pair{a}{b}$ and on $\Pair{b}{a}$, are independent (\cref{asm:ber_trial}).
Symmetry of $\ARelation{R}$ on this pair requires the two trials to \emph{agree}:
both report the pair, or both omit it. For a positive pair (both ordered tuples
lie in $\Relation{R}$) each trial is a true positive with probability $\tprate$
and a miss with probability $\fnrate$, so the two agree with probability
\begin{equation}
    \tprate^2 + \fnrate^2\,.
\end{equation}
For a negative pair each trial is a false positive with probability $\fprate$ and
a true negative with probability $1-\fprate$, so the two agree with probability
\begin{equation}
    (1-\fprate)^2 + \fprate^2\,.
\end{equation}
The pairs are disjoint in their trials, so the joint event is the product over all
$P$ positive and $N$ negative pairs. Each base is a sum of two squares of
probabilities summing to one, hence strictly less than one unless one probability
is zero; for interior rates the product vanishes as $P + N \to \infty$.
\end{proof}

Symmetry, unlike reflexivity, admits a cheap repair. Rather than \emph{hope} the
two ordered trials of every pair agree, we \emph{force} agreement by taking the
symmetric closure: admit a pair if either of its ordered trials fires. The next
theorem (the model worked result of this section) shows the closure restores the
invariant deterministically, at a quantified cost to the rates.

\begin{theorem}[Symmetric closure]
\label{thm:symmetric_closure}
	Let $\Relation{R} \subseteq \Set{A} \times \Set{A}$ be a \emph{symmetric} relation and let $\ARelation{R}[\fprate][\tprate]$ be a random approximate relation of $\Relation{R}$ with false-positive rate $\fprate$ and true-positive rate $\tprate$. Define the symmetric closure by
	\begin{equation}
		\Pair{x_1}{x_2} \in \ARelation{R}[s] \iff \Pair{x_1}{x_2} \in \ARelation{R} \;\text{ or }\; \Pair{x_2}{x_1} \in \ARelation{R}\,.
	\end{equation}
	Then $\ARelation{R}[s]$ is a random approximate symmetric relation
	\begin{equation*}
		\ARelation{R}[s]\!\left[1-(1-\fprate)^2\right]\!\left[1-(1-\tprate)^2\right]
	\end{equation*}
	with false-positive rate $1-(1-\fprate)^2$ and true-positive rate $1-(1-\tprate)^2$.
\end{theorem}
\begin{proof}
	Because $\Relation{R}$ is symmetric, $\Pair{x}{y} \in \Relation{R}$ if and only if $\Pair{y}{x} \in \Relation{R}$, so each unordered pair is tested through two trials, on $\Pair{x}{y}$ and on $\Pair{y}{x}$, which are independent by \cref{asm:ber_trial}. The closure admits the pair if \emph{either} trial fires.

	For a \emph{negative} pair ($\Pair{x}{y} \notin \Relation{R}$, hence $\Pair{y}{x} \notin \Relation{R}$ too), each trial is a false positive with probability $\fprate$, so the pair is admitted unless both trials miss:
	\begin{equation}
		\Prob{\Pair{x}{y} \in \ARelation{R}[s] \Given \Pair{x}{y} \notin \Relation{R}} = 1 - (1-\fprate)^2\,.
	\end{equation}
	For a \emph{positive} pair ($\Pair{x}{y} \in \Relation{R}$, hence $\Pair{y}{x} \in \Relation{R}$ too), each trial is a true positive with probability $\tprate$, so the pair is admitted unless both trials miss:
	\begin{equation}
		\Prob{\Pair{x}{y} \in \ARelation{R}[s] \Given \Pair{x}{y} \in \Relation{R}} = 1 - (1-\tprate)^2\,.
	\end{equation}
	The same doubling that raises the false-positive rate of the \textsc{or} also raises the true-positive rate; the true-positive slot is therefore $1-(1-\tprate)^2$, not the unchanged $\tprate$. For example, at $\tprate = 0.9$ the closure true-positive rate is $1 - (0.1)^2 = 0.99$.
\end{proof}

The closure is the template for repairing a fragile invariant: replace a
conjunction of hopeful trials by an operator that enforces the property outright,
and read off the rate it costs. We use this same closure freely when composing
binary relations that must carry the symmetric property.

\subsection{Single-valuedness}
\label{sec:functional}
A relation $\Relation{R} \subseteq \Set{X} \times \Set{Y}$ is \emph{functional}
(single-valued, a partial function) when the selection $\Selection_x \Relation{R}$
has cardinality $0$ or $1$ for every key $x \in \Set{X}$: each key carries at most
one value. The approximation can break this property in only one direction. A
false negative on the true value of a key leaves that key with zero values, which
is still at most single-valued; only a \emph{false positive}, which adds a
spurious second value, violates functionality. Single-valuedness is therefore
preserved exactly when no key acquires a spurious value.

\begin{theorem}[Single-valuedness is preserved with vanishing probability]
\label{thm:functional}
Let $\PFun{f} \subseteq \Set{X} \times \Set{Y}$ be functional. For a key $x$, the
approximation $\ARelation{f}[\fprate][\tprate]$ assigns no spurious value to $x$
with probability $(1-\fprate)^{\Card{\Set{Y}}-1}$, since each of the
$\Card{\Set{Y}}-1$ values other than the true one is a false positive with
probability $\fprate$ and these trials are independent (\cref{asm:ber_trial}). Over
all keys with a defined value, the probability that $\ARelation{f}$ remains
single-valued is
\begin{equation}
    \prod_{x} (1-\fprate)^{\Card{\Set{Y}}-1}
    = (1-\fprate)^{\Card{\Dom(\PFun{f})}\,(\Card{\Set{Y}}-1)}\,,
\end{equation}
which tends to $0$ as the value-universe $\Set{Y}$ or the domain grows, for any
$\fprate > 0$.
\end{theorem}
\begin{proof}
Fix a key $x \in \Dom(\PFun{f})$ with true value $y_x$. The pairs
$\SetBuilder{\Pair{x}{y}}{y \in \Set{Y}, y \neq y_x}$ are all negative for
$\PFun{f}$, so each is reported by $\ARelation{f}$ with probability $\fprate$,
independently. The key $x$ stays single-valued exactly when none of these
$\Card{\Set{Y}}-1$ false-positive trials fires, with probability
$(1-\fprate)^{\Card{\Set{Y}}-1}$. (Losing the true value $y_x$, an event of
probability $\fnrate$, leaves $x$ with no value, which does not violate
single-valuedness, so it does not enter the product.) Distinct keys index disjoint
pair sets, so the events are independent and the joint probability is the product
over the $\Card{\Dom(\PFun{f})}$ defined keys. Each factor is below one for
$\fprate > 0$, so the product vanishes as the exponent grows. This is the
functional reading of the approximate partial function of
\cref{def:approx_pfun}.
\end{proof}

\subsection{Transitivity}
\label{sec:transitive}
Transitivity is the hardest of the four to preserve and, unlike symmetry, has no
inexpensive closure. A relation is \emph{transitive} when
$\Pair{x}{y} \in \Relation{R}$ and $\Pair{y}{z} \in \Relation{R}$ together imply
$\Pair{x}{z} \in \Relation{R}$. The constraint ranges over \emph{triples}, not
pairs, and the implicated trials overlap (a single pair $\Pair{y}{z}$ participates
in many triples), so the preservation probability does not factor into independent
per-triple terms the way reflexivity and symmetry do.

We state the honest limiting behavior rather than a spurious closed form. For each
ordered triple $(x,y,z)$ whose two premises hold in $\Relation{R}$, transitivity
requires the conclusion $\Pair{x}{z}$ to be reported by $\ARelation{R}$. If
$\Pair{x}{z} \in \Relation{R}$ the conclusion holds with probability $\tprate$;
if $\Pair{x}{z} \notin \Relation{R}$ it holds (by a compensating false positive)
with probability $\fprate$. There are on the order of $\Card{\Set{A}}^3$ such
triple constraints, far more than the $\Card{\Set{A}}$ diagonal constraints of
reflexivity or the $\binom{\Card{\Set{A}}}{2}$ pair constraints of symmetry, and
the trials are coupled, so the preservation probability decays the fastest of the
four and tends to $0$ well before the others. The leading, most fragile behavior
is governed by the conclusion pairs that must be \emph{manufactured} by false
positives (those with $\Pair{x}{z} \notin \Relation{R}$): as a first-order
heuristic, ignoring the cross-triple coupling just noted, each such conclusion
contributes a factor near $\fprate \ll 1$, and the true rate is smaller still once
that coupling is accounted for.

A transitive-closure repair, the analogue of the symmetric closure, is far more
involved. Adding a manufactured pair $\Pair{x}{z}$ to satisfy one triple can pair
with existing pairs to create \emph{new} required conclusions, which may in turn
require further pairs, a cascade with no fixed per-pair rate cost. We therefore
leave transitivity as the honest limiting result: preservation probability tends
to zero fastest among the four invariants, and a deterministic repair, if desired,
must be computed as a genuine closure rather than read off as a doubled rate.

\subsection{A closure dichotomy}
\label{sec:closure_dichotomy}

The symmetric closure of \cref{thm:symmetric_closure} repaired a fragile invariant
at a fixed rate cost, while transitivity (\cref{sec:transitive}) admitted no
comparable repair. The contrast is not an accident of two examples; it is a
dichotomy, and the dividing line is whether the invariant's constraints
\emph{decompose into independent, bounded groups} or \emph{cascade}.

Call an invariant \emph{local} if the pairs it constrains partition into disjoint
bounded groups, the \emph{orbits}, of at most $m$ ordered pairs, with the
invariant holding for $\ARelation{R}$ exactly when every orbit is admissible on
its own. The monotone local invariants we treat are \emph{class-induced}: the
orbits are the classes of an equivalence on the tuple universe, and a relation
satisfies the invariant exactly when it is a \emph{union of orbits}. Symmetry is
class-induced with $m=2$: the orbit of $\{a,b\}$ is
$\{\Pair{a}{b}, \Pair{b}{a}\}$, and a relation is symmetric exactly when it
contains both elements of every orbit it meets. Reflexivity is the degenerate
\emph{forced} case $m=1$: each diagonal pair $\Pair{a}{a}$ is its own orbit,
whose value the invariant forces to true. Larger orbits arise as soon as the
invariant quotients by a larger symmetry: a $k$-ary relation required to be
invariant under cyclic shifts of its coordinates has orbits of size dividing
$k$, and one required to be invariant under all coordinate permutations (a
relation of multisets) has orbits of size up to $k!$. Two consequences of the
class-induced definition do the work below. First, because the exact
$\Relation{R}$ satisfies the invariant, $\Relation{R}$ is itself a union of
orbits, so every orbit is \emph{homogeneous}: wholly inside $\Relation{R}$ or
wholly outside it. Second, the constraints decompose orbit by orbit, so a repair
may treat orbits independently. A repair that adds pairs is \emph{monotone};
one that deletes pairs is \emph{antitone}. Single-valuedness is the antitone
case and is \emph{not} class-induced: the orbit of a key $x$ is its row
$\{\Pair{x}{y} : y \in \Set{Y}\}$, admissible when at most one pair is present,
a cardinality cap rather than a union-of-classes condition.

\begin{theorem}[Class-induced invariants admit a one-pass closure]
\label{thm:local_closure}
Let $P$ be a class-induced invariant with orbits of size at most $m$ (so that,
$\Relation{R}$ satisfying $P$, every orbit is homogeneous), repaired by
the \emph{OR-closure} that admits every pair of an
orbit as soon as any pair of that orbit is reported. The closure makes every orbit
internally consistent in one pass, so its output satisfies $P$ deterministically.
For $m \ge 2$ it combines the orbit's $m$ trials, yielding a random approximate
relation of false-positive rate $1-(1-\fprate)^{m}$ and
true-positive rate $1-(1-\tprate)^{m}$, a cost independent of $\Card{\Set{A}}$. The
degenerate orbit size $m=1$ (reflexivity) carries no second trial to combine: its
orbit is a single known required pair, restored exactly by inserting the fixed
diagonal at true-positive rate $1$ and no false positives, which is free.
\end{theorem}

\begin{proof}
Within an orbit the $m$ ordered trials are independent (\cref{asm:ber_trial}), and
the OR-closure admits the whole orbit unless all $m$ miss. For an orbit of true
non-members each trial is a false positive with probability $\fprate$, so the orbit
is admitted with probability $1-(1-\fprate)^{m}$; for an orbit of true members each
trial is a true positive with probability $\tprate$, so it is admitted with
probability $1-(1-\tprate)^{m}$. An admitted orbit is filled and hence internally
consistent, so the output satisfies $P$ on every orbit; the orbits are disjoint and
the OR-closure of an orbit adds only pairs of that same orbit, so the repairs do
not interfere and one pass suffices. For $m=2$ this recovers
\cref{thm:symmetric_closure}. The $m=1$ orbit has a single trial and no agreement
to enforce, so the OR-closure cannot lift its true-positive rate above $\tprate$;
reflexivity is instead secured by inserting the known diagonal outright, exactly
and at no rate cost. Homogeneity is what makes the two rates clean; were an
orbit permitted to straddle $\Relation{R}$, its fill probability would be the
heterogeneous product $1-(1-\tprate)^{j}(1-\fprate)^{m-j}$ in the member count
$j$, the same form as the mixed-fiber projection rate, but a class-induced
invariant never produces such an orbit.
\end{proof}

The OR-closure restores the invariant by spending: it consults all $m$ trials of
an orbit and moves both rates. It is natural to ask what the invariant alone
costs, and the answer is nothing.

\begin{theorem}[The canonical test: the invariant itself is free]
\label{thm:canonical_test}
Let $P$ be class-induced with orbits of size at most $m$. Fix a representative
$\operatorname{rep}(O)$ in each orbit $O$, and let the \emph{canonical test} report
$t \in \ARelation{R}[c]$ exactly when the single trial on
$\operatorname{rep}(\operatorname{orbit}(t))$ fires. Then $\ARelation{R}[c]$ satisfies $P$
deterministically, with false-positive rate $\fprate$ and true-positive rate
$\tprate$ unchanged, and a structure built for the canonical test stores one
element per member orbit, $\Card{\Relation{R}}/2$ elements for a symmetric
relation with empty diagonal, rather than $\Card{\Relation{R}}$. Forced orbits
(the reflexive diagonal) need no trial at all: answer true, exactly. The
invariant itself therefore costs nothing, in rate or in space.
\end{theorem}

\begin{proof}
Membership under the canonical test depends only on $\operatorname{orbit}(t)$, so
$\ARelation{R}[c]$ is a union of orbits and satisfies $P$ in every realization.
By homogeneity, $\operatorname{rep}(\operatorname{orbit}(t)) \in \Relation{R}$ if and only if
$t \in \Relation{R}$, so the canonical trial is a true-positive trial (rate
$\tprate$) precisely on members and a false-positive trial (rate $\fprate$)
precisely on non-members. Only representatives are ever consulted, so a
purpose-built structure stores one element per member orbit. The same test
applies post hoc to an approximation built without regard to orbits, by
rewriting each query to its representative: the invariant is restored with no
rebuild and no rate movement.
\end{proof}

The two repairs now sit side by side. The canonical test secures the invariant
at base rates and base space, and is available both at construction time and,
by query rewriting, post hoc. The OR-closure spends more to get more: it
consults all $m$ trials of the orbit and lifts the true-positive rate from
$\tprate$ to $1-(1-\tprate)^{m}$, at the false-positive cost
$1-(1-\fprate)^{m}$. The surcharge computed next is therefore the price of the
\emph{recall boost}, not of the invariant.

The single-valued (antitone) case is repaired in one pass the other way, by
\emph{deletion}: keep at most one reported value per key. It too is local and
$\Card{\Set{A}}$-independent, trading a lowered false-positive rate against the
risk of dropping the true value, the selection counterpart of \cref{thm:functional}.
The rate is closed-form under uniform tie-breaking: a defined key keeps its true
value exactly when the value is reported (probability $\tprate$) and survives a
uniform draw among the $1+K$ reported candidates, where $K$ counts the spurious
reports, $K \sim \operatorname{Binomial}(\Card{\Set{Y}}-1, \fprate)$; averaging
$1/(1+K)$ over the binomial gives the kept-value rate
\begin{equation}
  \tprate \, \frac{1-(1-\fprate)^{\Card{\Set{Y}}}}{\fprate \, \Card{\Set{Y}}}\,,
\end{equation}
which tends to $\tprate$ as $\fprate \to 0$ and decays like
$\tprate/(\fprate \Card{\Set{Y}})$ once $\fprate \Card{\Set{Y}} \gg 1$. Securing
single-valuedness \emph{by construction} is instead a change of type: store the
partial function as a map, key to value, rather than as a relation, the object
of the maps companion~\cite{bernoulliMaps}.

The monotone cost has a clean reading in the space-accuracy duality of the entropy
companion~\cite{bernoulliEntropy}, where a positive approximate set spends
$-\log_2 \fprate$ bits per element.

\begin{corollary}[The space cost of a local invariant]
\label{cor:invariant_space_cost}
It suffices to run the underlying structure at base rate $\fprate =
1-(1-\fprate')^{1/m} \approx \fprate'/m$ for the OR-closure of a monotone local
invariant of orbit size $m$ to hold at a target false-positive rate $\fprate'$,
which by the space-accuracy duality costs
\begin{equation}
  -\log_2 \fprate \;=\; -\log_2 \fprate' + \log_2 m + o(1)
\end{equation}
bits per element: an additive surcharge of $\log_2 m$ bits per element buys the
OR-closure's recall boost at an unmoved false-positive target. Symmetry's boost
costs one extra bit per element; reflexivity is free outright. The surcharge
prices the boost, not the invariant: \cref{thm:canonical_test} settles in the
negative the question of whether the invariant itself forces any surcharge,
since the canonical test preserves it at base rate and base space. The
OR-closure remains the right tool when the $m$ orbit trials already exist and
recall is worth buying.
\end{corollary}

Transitivity sits on the other side of the line, and the obstruction is exactly
the failure of decomposability. It fails in the strongest sense: not only does
the OR-closure not apply, no canonical test can exist at all.

\begin{lemma}[No equivalence induces transitivity]
\label{lem:no_quotient}
For $\Card{\Set{A}} \ge 2$, no equivalence on $\Set{A} \times \Set{A}$ has the
transitive relations as its unions of classes. Transitivity therefore admits no
canonical test, and the cascade below is not an artifact of choosing a poor
repair.
\end{lemma}

\begin{proof}
Unions of classes of a fixed equivalence are closed under arbitrary union. The
transitive relations are not: $\{\Pair{a}{b}\}$ and $\{\Pair{b}{a}\}$ are each
transitive, and their union omits $\Pair{a}{a}$ and so is not.
\end{proof}

\begin{theorem}[Transitivity cascades]
\label{thm:transitivity_cascade}
The transitive closure is not a one-pass repair. On the path relation
$\Relation{R} = \{\Pair{1}{2}, \Pair{2}{3}, \ldots, \Pair{k}{k+1}\}$ over
$\Card{\Set{A}} = k+1$ elements, reaching the transitive closure takes
$\lceil \log_2 k \rceil$ successive \emph{simultaneous} passes, each pass adding
at once every pair the current relation implies,
$\Relation{R} \gets \Relation{R} \cup (\Relation{R} \circ \Relation{R})$, and
manufactures $\binom{k+1}{2} =
\Theta(k^2)$ pairs. Both grow with $\Card{\Set{A}}$, so no fixed-cost,
$\Card{\Set{A}}$-independent repair of the form of \cref{thm:local_closure} exists.
\end{theorem}

\begin{proof}
One pass adds $\Pair{x}{z}$ whenever $\Pair{x}{y}$ and $\Pair{y}{z}$ are present,
producing on the path the distance-two pairs $\Pair{i}{i+2}$. A distance-four pair
needs a present distance-two pair, so it appears only on the second pass, and in
general distance-$2^{j}$ pairs appear on pass $j$; the longest distance $k$ is
reached after $\lceil \log_2 k \rceil$ passes. The closure is the full strict order
$\{\Pair{i}{j} : i < j\}$, of size $\binom{k+1}{2}$. The triples of transitivity
share pairs and chain, so the constraints cannot be partitioned into independent
orbits, which is the hypothesis \cref{thm:local_closure} requires.
\end{proof}

The four invariants therefore split on one line. Reflexivity, symmetry, and
single-valuedness are \emph{local}: their constraints decompose into independent
bounded orbits, so each is secured in a single pass at a cost independent of the
universe, for free by the canonical test where the invariant is class-induced
(\cref{thm:canonical_test}), at a $\log_2 m$ surcharge where the OR-closure's
recall boost is wanted (\cref{cor:invariant_space_cost}), and by deletion at a
closed-form rate for the antitone one. Transitivity
is \emph{global}: no equivalence induces it (\cref{lem:no_quotient}), its
triples chain, the repair cascades to depth
$\lceil \log_2 k \rceil$, and no bounded closure exists. The cut is decomposability,
not the superficial pair-versus-triple arity.

That some integrity constraints are far cheaper to repair than others is a
familiar theme in the database-repair and consistent-query-answering
tradition~\cite{bertossi-2011-repairing}, where the complexity of repairing or
querying an inconsistent database turns on the constraint class. The dichotomy
here runs on a different axis. It is not the computational complexity of a repair
over an exact database, but the closed-form \emph{rate and space cost} of
restoring an invariant in the representation-error setting, where the relation is
a compressed approximate structure and the repair is a one-pass canonical test or
closure whose price is zero (\cref{thm:canonical_test}) or the $\log_2 m$ recall
surcharge of \cref{cor:invariant_space_cost}.

These four results share a model-order reading. Each invariant partitions the
universe into pairs (or triples) that play different error roles (diagonal versus
off-diagonal; positive versus negative; premise versus conclusion), so an
approximation forced to honor the invariant is a higher-order Bernoulli relation
whose blocks carry distinct rates, exactly the model-order versus parameter-count
distinction of the underlying set model~\cite{bernoulliSets}.

\begin{remark}
Some of these properties may be difficult to guarantee, especially for cipher relations.
However, \emph{tractability} is a separate issue from specification and definition.
\end{remark}

\section{The approximate partial function}
\label{sec:map}
A \emph{partial function} is the single-valued special case of a relation: each tuple in the domain maps to at most one tuple in the range. Approximate-map composition, the error rates that result when two approximate partial functions are composed, is treated in the maps paper~\cite{bernoulliMaps}, whose composition theorem gives a false-negative rate $\fnrate_{g \circ f} = 1 - (1-\fnrate_1)(1-\fnrate_2)$ with an explicit overlap parameter governing the resulting false-positive rate. We do not duplicate that analysis here; what follows is the partial-function abstract data type viewed as a relation.

%%% Phase A (2026-06-08): the "Random approximate maps" machinery that stood here (map
%%% notation; Dom/Codom/Dod/Range; augmented codomain; the map-composition theorem with
%%% true-positive rate tau1*tau2 and the undefined IntervalSpan false-positive rate; the
%%% pointwise-operations stub) was cut. That material is owned, and proved correctly, by
%%% the maps paper (\cite{bernoulliMaps}); the composition theorem there gives
%%% omega_{g.f} = 1 - (1-omega1)(1-omega2) with an explicit overlap parameter for the FPR.
%%% Only the genuinely-relational partial-function ADT (def:approx_pfun, below) is kept.

We consider relations known as \emph{partial functions} as given by the following definition.

\begin{definition}
	A relation
	\begin{equation}
	\PFun{f} \subseteq \Set{X}[1] \times \cdots \times \Set{X}[n] \times \Set{Y}[1] \times \cdots \times \Set{Y}[m]\,,
	\end{equation}
	is a \emph{function} if $\Card{\Proj_{\Set{X}[1] \times \cdots \times \Set{X}[n]}(\PFun{f})} = \Card{\PFun{f}}$ and is a \emph{partial} function if, additionally, $\Card{\Proj_{\Set{X}[1] \times \cdots \times \Set{X}[n]}(\PFun{f})} < \Card{\Set{X}[1] \times \cdots \times \Set{X}[n]}$.
	
	We say that each tuple $\Tuple{x_1,\ldots,x_n}$ in the relation \emph{maps} to a tuple $\Tuple{y_1,\ldots,y_m} \in \Set{Y}[1] \times \cdots \times \Set{Y}[m]$. To make the mapping explicit, we denote the partial function by
	\begin{equation}
	\PFun{f} \colon \Set{X}[1] \times \cdots \times \Set{X}[n] \mapsto \Set{Y}[1] \times \cdots \times \Set{Y}[m]\,.
	\end{equation}
\end{definition}
A partial function represents \emph{many}-to-\emph{one} relationships, so a tuple in its domain unambiguously names its mapped value. The projection $\Proj_{\Set{X}[1] \times \cdots \times \Set{X}[n]}(\PFun{f})$ is the \emph{domain} of $\PFun{f}$, denoted $\Dom(\PFun{f})$; the projection $\Proj_{\Set{Y}[1] \times \cdots \times \Set{Y}[m]}(\PFun{f})$ is its \emph{image}, $\Image(\PFun{f})$; the \emph{range} is $\Range(\PFun{f}) = \Set{Y}[1] \times \cdots \times \Set{Y}[m]$. \emph{Maps} and \emph{partial functions} are synonymous; \emph{map} is the more common term when the domain elements are called \emph{keys} and the range elements \emph{values}. The classical development (counting, table representation, worked examples) belongs to the maps companion~\cite{bernoulliMaps}; here we need only the abstract data type of the approximate partial function, which the single-valuedness results of \cref{sec:functional,sec:closure_dichotomy} ground.
\begin{definition}
	\label{def:approx_pfun}
	A relation $\APFun{f} = \Set{X}[1] \times \cdots \times \Set{X}[n] \times \Set{Y}[1] \times \cdots \times \Set{Y}[m]$ is an \emph{approximate partial function} of $\PFun{f} \colon \Set{X}[1] \times \cdots \times \Set{X}[n] \mapsto \Set{Y}[1] \times \cdots \times \Set{Y}[m]$ if the following conditions hold:
	
	Let a random tuple that is selected uniformly at random from the universe $\Set{X}[1] \times \cdots \times \Set{X}[n]$ be denoted by $\Tuple{\RV{X}} = \Tuple{\RV{X_1},\cdots,\RV{X_n}}$. The partial function $\APFun{f}$ is an approximate partial function of $\PFun{f}$ with a false domain rate $\fprate$ and false complementary domain rate $\fnrate$ if the following conditions hold:
	\begin{enumerate}[(i)]
		\item If $\Tuple{\RV{X}}$ is in the domain of $\PFun{f}$, it is not in the domain of $\APFun{f}$ with a probability $\fnrate$,
		\begin{equation}
		\Prob{\Tuple{\RV{X}} \notin \Dom(\APFun{f}) \Given \Tuple{\RV{X}} \in \Dom(\PFun{f})} = \fnrate\,.
		\end{equation}    
		\item If $\Tuple{\RV{X}}$ is \emph{not} in the domain of $\PFun{f}$, it is in the domain of $\APFun{f}$ with a probability $\fprate$,
		\begin{equation}
		\Prob{\Tuple{\RV{X}} \in \Dom(\APFun{f}) \Given \Tuple{\RV{X}} \notin \Dom(\PFun{f})} = \fprate\,.
		\end{equation}
	\end{enumerate}
\end{definition}

\end{document}